% 八物品教学目录；节点得分仅用于TDM束宽为2的演示。
\begin{tikzpicture}[
 font=\figurefont\footnotesize,
 branch/.style={rectangle,draw=BookRule,fill=BookPaper,
   align=center,minimum width=12mm,minimum height=8mm,inner sep=3pt},
 kept/.style={branch,draw=FigureModel,fill=FigureModel!10},
 leaf/.style={branch,minimum width=8mm,minimum height=7mm},
 edge/.style={draw=BookMuted,line width=.8pt},
 active/.style={draw=FigureModel,line width=1.1pt},
 cut/.style={draw=BookMuted,dashed,line width=.8pt}
]
 \node[branch] (root) at (0,0) {目录};
 \node[kept] (a) at (-2.6,-1.1) {$A$};
 \node[kept] (b) at (2.6,-1.1) {$B$};
 \node[kept] (a1) at (-3.9,-2.3) {$A_1$\\0.9};
 \node[branch] (a2) at (-1.3,-2.3) {$A_2$\\0.6};
 \node[kept] (b1) at (1.3,-2.3) {$B_1$\\0.8};
 \node[branch] (b2) at (3.9,-2.3) {$B_2$\\0.3};
 \node[leaf,fill=FigureOutput!10] (i1) at (-4.55,-3.6) {$a$};
 \node[leaf,fill=FigureOutput!10] (i2) at (-3.25,-3.6) {$b$};
 \node[leaf] (i3) at (-1.95,-3.6) {$c$};
 \node[leaf] (i4) at (-.65,-3.6) {$d$};
 \node[leaf,fill=FigureOutput!10] (i5) at (.65,-3.6) {$e$};
 \node[leaf,fill=FigureOutput!10] (i6) at (1.95,-3.6) {$f$};
 \node[leaf] (i7) at (3.25,-3.6) {$g$};
 \node[leaf] (i8) at (4.55,-3.6) {$h$};
 \draw[active] (root) -- (a);
 \draw[active] (root) -- (b);
 \draw[active] (a) -- (a1);
 \draw[edge] (a) -- (a2);
 \draw[active] (b) -- (b1);
 \draw[edge] (b) -- (b2);
 \draw[active] (a1) -- (i1);
 \draw[active] (a1) -- (i2);
 \draw[active] (b1) -- (i5);
 \draw[active] (b1) -- (i6);
 \draw[cut] (a2) -- (i3);
 \draw[cut] (a2) -- (i4);
 \draw[cut] (b2) -- (i7);
 \draw[cut] (b2) -- (i8);
 \node[align=center,text=BookMuted] at (0,-4.4)
   {第二层保留$A_1,B_1$；虚线路径不再展开};
\end{tikzpicture}
